\chapter{Concept Testing}

\section{Session Overview}

Concept testing is an early-stage research activity used to gather feedback on a proposed concept before committing to full development. It is different from product testing: concept testing asks whether the idea is attractive, understandable, and worth pursuing, while product testing checks the performance of a developed product.

By the end of this session, you should be able to:
\begin{enumerate}
    \item define what your concept test needs to learn;
    \item identify an appropriate survey population and sample size approach;
    \item choose a survey format that fits the concept and audience; and
    \item decide how to communicate the selected concept clearly to respondents.
\end{enumerate}

\section{Key Terms}

\begin{itemize}
    \item \textbf{Concept testing}: obtaining user feedback on an idea before detailed development is completed.
    \item \textbf{Survey population}: the group of people from whom feedback will be collected.
    \item \textbf{Market segment}: a category of customers sharing similar needs or contexts.
    \item \textbf{Sample size}: the number of respondents included in the study.
    \item \textbf{Survey format}: the method used to collect feedback, such as interview, online form, or email survey.
\end{itemize}

\section{Activity 1: Define the Purpose of the Concept Test}

Start by deciding which decisions the concept test should support. Typical concept-testing questions include:
\begin{itemize}
    \item Which concept do users prefer?
    \item What part of the concept is unclear or unattractive?
    \item What should be improved before further development?
    \item Does the concept appear useful enough to justify continuing the project?
\end{itemize}

For a hydroponic concept, an example research question could be: ``Would hobby growers trust and value an add-on product that provides low-cost nutrient-level alerts?''

\section{Activity 2: Choose the Survey Population}

Use the target market defined earlier in the project to identify a realistic survey population. For example, your concept may target home growers, educational users, or commercial hydroponic operators.

When selecting respondents, consider:
\begin{itemize}
    \item whether they resemble the intended users;
    \item whether they have enough experience to evaluate the concept meaningfully; and
    \item whether the selected group reflects the concept's price level, usage context, and maintenance expectations.
\end{itemize}

If you estimate sample size using an online calculator, record the assumptions used, such as confidence level and margin of error. For classroom use, representativeness and justification are more important than achieving a very large sample.

\section{Activity 3: Choose the Survey Format and Draft Questions}

Select a survey format that matches the target users and available resources. Common options include:
\begin{itemize}
    \item face-to-face interviews;
    \item telephone interviews;
    \item email surveys; and
    \item online forms.
\end{itemize}

A simple concept test usually includes:
\begin{enumerate}
    \item a short introduction;
    \item screening questions to confirm the respondent fits the target market;
    \item presentation of the concept;
    \item evaluation questions; and
    \item closing or demographic questions if needed.
\end{enumerate}

\begin{table}[htbp]
\centering
\caption{Example concept-test planning template}
\label{tab:session12-test-plan}
\renewcommand{\arraystretch}{1.2}
\begin{tabularx}{\textwidth}{|>{\raggedright\arraybackslash}p{0.24\textwidth}|X|}
\hline
\textbf{Plan element} & \textbf{Draft content} \\
\hline
Purpose of test &  \\
\hline
Target market segment &  \\
\hline
Survey population &  \\
\hline
Sample size approach &  \\
\hline
Survey format &  \\
\hline
Main evaluation questions &  \\
\hline
Communication method &  \\
\hline
\end{tabularx}
\end{table}

\subsection{Sample Evaluation Questions}

The following question types are commonly useful:
\begin{itemize}
    \item How clear is the concept description?
    \item How useful would this product be in your hydroponic setup?
    \item What concerns would you have about using this concept?
    \item How likely would you be to try or purchase this product if available?
    \item What feature would most improve the concept for you?
\end{itemize}

\section{Activity 4: Decide How to Communicate the Concept}

Your concept must be presented clearly enough for respondents to evaluate it. Use the simplest communication method that still explains the important features.

\begin{table}[htbp]
\centering
\caption{Communication method comparison for concept testing}
\label{tab:session12-communication-matrix}
\renewcommand{\arraystretch}{1.2}
\begin{tabularx}{\textwidth}{|>{\raggedright\arraybackslash}p{0.22\textwidth}|>{\raggedright\arraybackslash}p{0.18\textwidth}|>{\raggedright\arraybackslash}p{0.18\textwidth}|X|}
\hline
\textbf{Method} & \textbf{Fidelity} & \textbf{Effort} & \textbf{When it is useful} \\
\hline
Verbal description & Low & Low & Early-stage concepts with simple value propositions \\
\hline
Sketch or storyboard & Medium & Low to medium & Concepts needing visual explanation of layout or workflow \\
\hline
Photo, rendering, or slide & Medium & Medium & Concepts where appearance and arrangement matter \\
\hline
Video or simulation & High & Medium to high & Concepts involving interaction, movement, or system behaviour \\
\hline
Physical model or prototype & High & High & Concepts requiring tactile or spatial understanding \\
\hline
\end{tabularx}
\end{table}

For example, a hydroponic add-on that changes nutrient flow automatically may be easier to explain through a storyboard or short video than through text alone.

\section{In-Class Consolidation}

In class, your team should compare the individual plans and agree on:
\begin{enumerate}
    \item the final purpose of the concept test;
    \item the target respondents and sample-size approach;
    \item the survey format;
    \item the draft survey questions; and
    \item the communication method used to show the concept.
\end{enumerate}

Before finalising the survey, check the wording for bias, clarity, and relevance. If possible, pilot the plan with a small number of peers and revise unclear questions.

\section{Ethics and Documentation}

When collecting feedback, explain the purpose of the survey, treat respondents respectfully, and avoid collecting unnecessary personal information. Keep a record of the final plan, the questions used, and the rationale for your choices.

\section{Summary and Next Steps}

At the end of this session, your group should have a documented concept-test plan that explains what you want to learn, who you will ask, how you will ask them, and how you will present the concept. The results of that test should inform concept refinement and the final project narrative.
